\SourceLeaf{278}
\par\indent Die Zahlenwerte in unserem Beispiel (für den Logarithmus nach Briggs und bei Angabe von \(q\) in Grad) sind
\[
\begin{aligned}
\log m={}&-0{,}0049796163\ 94(\frac{q}{100})^3\\
&-0{,}0016150307\ 6(\frac{q}{100})^4\\
&-0{,}0023973954(\frac{q}{100})^5\\
&-0{,}0125671(\frac{q}{100})^6
\end{aligned}
\]
\Article{10}
\indent Bei einer umfangreichen Vermessung, bei der sehr viele Punkte vom Sphäroid auf die Kugel oder umgekehrt übertragen werden, wird man lieber ein für alle Mal eine umfangreiche Tafel berechnen, statt jedes Mal auf die Formeln zurückzugreifen. Eine solche Tafel lässt sich aber bequemer verwenden, wenn man die Breite auf der Kugel, \(Q+q\), als Argument wählt, statt die Breite auf dem Sphäroid zu verwenden. Der Übergang von dieser zur Breite auf dem Sphäroid ist nämlich viel häufiger erforderlich als der umgekehrte. Wer im Rechnen erfahren ist, braucht wohl keinen Hinweis darauf, dass man für eine solche Tafel nur eine mäßige Anzahl von Einträgen direkt berechnet. Aus ihnen lassen sich die übrigen durch ein geeignetes Interpolationsverfahren mit derselben Genauigkeit und sehr wenig Mühe bestimmen. Dazu kommen also die in Artikel 8 und 9 angegebenen Reihen zur Anwendung. Gerade deshalb ist es vorteilhaft, bei \(P\) auf eine runde Zahl zu verzichten und für \(Q\) eine zu wählen.

Am Schluss dieser Abhandlung füge ich eine solche Tafel bei. Ihr liegt der Normalwert \(Q=52^0 40'\) zugrunde (wie dem bisher betrachteten Beispiel), und sie erstreckt sich über zwölf Grad, von \(46^0 40'\) bis \(58^0 40'\), für alle Werte des Arguments \(Q+q\) in Schritten von einer Minute. Sie gibt den zugehörigen Wert von \(P+p\) auf fünf Dezimalstellen der Sekunde genau und den Logarithmus nach Briggs von \(m\) auf zehn Stellen, und zwar in Einheiten der zehnten Dezimalstelle. Außerdem gibt sie den in Sekunden ausgedrückten Wert von \(-\tfrac{\mathrm{d}m}{2m\mathrm{d}q}\) an. Die Verwendung dieser letzten Spalte wird weiter unten erklärt. Ich habe die Tafel mit so vielen Dezimalstellen angegeben, damit sie auch für die genaueste Berechnung einer trigonometrischen Vermessung, nämlich für eine
